\begin{frame}[noframenumbering,plain]
    \setcounter{framenumber}{1}
    \maketitle
\end{frame}
\begin{frame}
    \frametitle{Введение}

\end{frame}


\begin{frame}
    \frametitle{Положения, выносимые на защиту}
    \begin{itemize}
    \item Сформулирована адаптивная стоимостная модель с процедурой онлайн калибровки параметров по телеметрии реальных запросов и состоянию буферного кэша; экспериментально установлено их улучшение по сравнению с базовой конфигурацией оптимизатора.

    \item Введено представление плана выполнения в эмбеддинговом пространстве с помощью большой языковой модели и предложен алгоритм переноса подсказок оптимизатору на основе близости план–к–плану с двухэтапной проверкой согласованности/безопасности; показано, что при заданных критериях согласованности и безопасности метод устойчиво улучшает время выполнения относительно базовой конфигурации.

    \item Предложен агентно-поисковый подход с участием большой языковой модели в замкнутом контуре обратной связи с СУБД; формализованы компоненты контура и установлено, что при фиксированном бюджете итераций агент достигает улучшений, сравнимых или превосходящих одношаговые/жадные стратегии подбора подсказок.
    \end{itemize}
\end{frame}

\note{
    Проговариваются вслух положения, выносимые на защиту
}

\begin{frame}
    \frametitle{Содержание}
    \tableofcontents
\end{frame}

